\chapter*{Til lesaren}
\addcontentsline{toc}{chapter}{Til lesaren}
\markboth{Til lesaren}{}

\begin{center}\itshape
Dette er motboka. Ho er skriven for å vinne.
\end{center}
\vspace{8mm}

Det finst ei anna bok i dette prosjektet, \textsc{Grunnlaust}, som hevdar at designfaget aldri nådde eit haldbart fundament — at det manglar ei felles etikk, ein delt metode, eit presist fagspråk og ei eiga stemme, og at det difor er ein avleggar i vente. Denne boka tek det motsette synet. Ikkje som ei pliktøving, og ikkje som ein stråmann sett opp for å bli velta. Eg har prøvt å gje forsvaret for designfaget den sterkaste forma eg er i stand til — å lese faget sine fremste tenkjarar, og særleg dei samtidige og dei norske, med all den velvilja dei har krav på, og å byggje deira beste argument til ein heilskap som faktisk kan vinne diskusjonen.

Grunnen til at begge bøkene finst, er enkel. Eit syn som ikkje toler sitt eige sterkaste motsvar, fortener ikkje å bli ståande. \textsc{Grunnlaust} sin diagnose er hard, og om han er rett, må han kunne stå imot det beste forsvaret faget kan mønstre. Så her er det forsvaret, samla og skjerpa. Lesaren som har lese begge, vil sjølv måtte avgjere kven som har rett — og det er heile poenget. Eg skriv ikkje for å fortelje deg kva du skal meine. Eg skriv for å gjere usemja så skarp som ho kan bli.

\section{Kva forsvaret faktisk hevdar}

Lat meg seie med ein gong kva denne boka \emph{ikkje} hevdar, for at forsvaret ikkje skal forvekslast med eit svakare eitt. Ho hevdar ikkje at designfaget er feilfritt, at det aldri har gjort skade, eller at alt er vel i utdanninga. Ho innrømmer mykje av det \textsc{Grunnlaust} skildrar som faktiske tilhøve. Det ho bestrider, er \emph{tolkinga} — slutninga frå desse tilhøva til at faget manglar eit fundament og difor ikkje fortener plassen sin.

Forsvaret kviler på fire påstandar, som boka byggjer ut ein etter ein. Den fyrste er at sjølve kravet er feilstilt: å krevje at design skal grunngje seg slik fysikken gjer, er ein kategorifeil, ikkje eit krav om strenghet. Den andre er at design \emph{har} ei eiga kunnskapsform — «designerly ways of knowing» — som er like reell og overførbar som kunnskapsformene i andre profesjonsfag, og som det finst ei førti år gammal forskingslitteratur om. Den tredje er at design har metode og kumulativ kunnskap, berre i eit anna register enn naturvitskapen: mønster, prinsipp, retningslinjer, og ein veksande forskingstradisjon. Den fjerde er at design har ei etikk og ei stemme — kodeksar, profesjonsorgan, ein politisk og samfunnsmessig rolle — som \textsc{Grunnlaust} oversér fordi dei ikkje har forma faget krev.

\section{Metoden: å lese med velvilje}

Eg har bygd boka på det same som \textsc{Grunnlaust}: på litteraturen. Men der \textsc{Grunnlaust} les faget mot håra, les denne boka det med håra. Eg har gått til dei samtidige designforskarane — internasjonalt og særleg i Noreg og Norden, der eit eige forskingsmiljø har vakse fram kring \textsc{Aho}, kring tidsskriftet \textsc{Formakademisk}, kring den systemorienterte designtradisjonen og den nordiske designforskinga — og spurt: kva er det \emph{beste} dei har å seie? Kvar gong eit argument kunne lesast både svakt og sterkt, har eg valt den sterke lesinga. Det er det ein stålmann er: motpartens posisjon i si beste form, ikkje i si verste.

Lesaren bør vite at forfattaren av denne boka òg er forfattaren av \textsc{Grunnlaust}, og at eg ikkje skjuler kvar mi eiga tvil ligg. Men i desse sidene er det forsvaret som fører ordet, og eg har gjort mitt for at det skal føre det godt. Døm sjølv om det held.

\section{Korleis boka er bygd}

Boka speglar \textsc{Grunnlaust} med vilje. Sju delar fører forsvaret fram: fyrst at spørsmålet er feilstilt (\textsc{del i}), så at design har ei kunnskapsform (\textsc{ii}), ein metode og kumulasjon (\textsc{iii}), at det vrange er ei innsikt og ikkje eit alibi (\textsc{iv}), at meining, menneske og det gode gjev faget grunn (\textsc{v}), og at dei samtidige og norske svara styrkjer forsvaret (\textsc{vi}). Til sist eit oppgjer (\textsc{vii}) der eg tek \textsc{Grunnlaust} sine påstandar éin for éin og svarar på dei i deira eigen sterkaste form. Kvart kapittel endar med \emph{Vidare lesing}.

Les denne boka saman med den andre. Kvar for seg er dei to einsidige med vilje. Saman er dei eit forsøk på å la den viktigaste usemja om designfaget bli ført så reieleg som eg klarte.

\vidarelesing{%
Kjernen i forsvaret er \textcite{cross2006} og \textcite{cross1982}; for design som eigen kunnskapskultur, sjå \textcite{schon1983}; for det integrerande synet, \textcite{buchanan1992wicked}. Motboka er \textsc{Grunnlaust} (same prosjekt), som denne boka svarar på.}
